\pdfoutput=1
\documentclass[11pt]{article}
\usepackage[margin=1in]{geometry}
\usepackage{amsmath,amssymb,amsthm}
\usepackage{enumitem}
\usepackage{hyperref}
\hypersetup{hidelinks,
  pdftitle={An explicit initial Theta-datum with mixed reduction at a split prime},
  pdfauthor={Toshihisa Urahigashi}}

\theoremstyle{plain}
\newtheorem{theorem}{Theorem}
\newtheorem{proposition}[theorem]{Proposition}
\newtheorem{corollary}[theorem]{Corollary}
\theoremstyle{definition}
\newtheorem{hypothesis}{Hypothesis}
\newtheorem{remark}{Remark}
\newtheorem{question}{Question}

\newcommand{\Fmod}{F_{\mathrm{mod}}}
\newcommand{\Vmod}{\mathbb{V}_{\mathrm{mod}}}
\newcommand{\Vbad}{\mathbb{V}^{\mathrm{bad}}}
\newcommand{\Vgood}{\mathbb{V}^{\mathrm{good}}}
\newcommand{\Vgoodmod}{\mathbb{V}^{\mathrm{good}}_{\mathrm{mod}}}
\newcommand{\VV}{\mathbb{V}}
\newcommand{\Ind}[1]{\textup{(Ind#1)}}
\newcommand{\SA}{\textup{(SA)}}
\newcommand{\ord}{\operatorname{ord}}
\newcommand{\Norm}{\operatorname{N}}

\title{An explicit initial $\Theta$-datum with mixed reduction\\
at a split prime}
\author{Toshihisa Urahigashi\\
\small\texttt{123026.urahigashi.toshihisa@gmail.com}}
\date{Version 0.1.5 --- 8 September 2026\\[3pt]
\small DOI \href{https://doi.org/10.5281/zenodo.22537342}{\texttt{10.5281/zenodo.22537342}} (all versions)\qquad ORCID \href{https://orcid.org/0009-0004-0460-6242}{\texttt{0009-0004-0460-6242}}}

\begin{document}
\maketitle

\begin{abstract}
We exhibit an explicit initial $\Theta$-datum in the sense of
\cite[Def.~3.1]{IUT1}: an elliptic curve in Legendre form over
$\mathbb{Q}(\sqrt5)$ with $\ell=7$, for which the rational prime $211$
\emph{splits} in the field of moduli, so that one place of multiplicative and one
place of good reduction lie above it.  Every condition (a)--(f) of
\cite[Def.~3.1]{IUT1} is verified individually, several by exact computation, and
the local invariants are $e_{\mathrm{g}}=1$, $e_{\mathrm{b}}=105$ and
$\ord_{211}(\underline{\underline q}_{\mathrm{b}})=29/7$, the profile being tame
at $211$ since $105\le209$.

The configuration is not incidental.  Because \cite[Def.~3.1(e)]{IUT1} makes
$\VV$ a \emph{section} of $\VV(K)\twoheadrightarrow\Vmod$, the tuples indexing a
local tensor packet at a rational prime are constant unless that prime has
several places in $\Vmod$; and for the curve constructed here the only odd
non-split places of multiplicative reduction are $3$ and $5$, carrying less than
$1.24\%$ of the relevant Tate contribution outside $2\ell$.  Nor can the
configuration be removed by a different choice of $\Vbad_{\mathrm{mod}}$: the
additional good-reduction hypothesis of \cite[Thm.~1.10]{IUT4} does not permit an
actual bad place outside $2\ell$ to be treated as chosen-good.

Under one hypothesis \SA{} about the source, isolated in \S\ref{s:sa} and
discussed in \S\ref{s:fals}, the upper bound obtained in
\cite[Thm.~1.10, proof, Step~(v)]{IUT4} by applying \cite[Prop.~1.2(iii)]{IUT4} to
a fixed target component is incompatible with the region it is applied to, with
margin $8/7$ at $j=1$ and $88/7$ at $j=2$.  We state \SA{} separately, list the
ways it can fail, and undertake to withdraw that application if the reading is
shown to be wrong; we would regard such a demonstration as the more valuable
outcome, since it would make explicit what the collection of possible images is
at a tuple of mixed reduction type.

We ask how the admitted configuration is treated (Question~\ref{q:hyp}).
The exact computations and the separate power-freeness certificate are
described in Section~\ref{s:repro}.
\end{abstract}

\medskip
\noindent\textbf{2020 Mathematics Subject Classification.}
Primary 14H25; Secondary 11G05, 14G32, 11R58.

\noindent\textbf{Keywords.}
inter-universal Teichm\"uller theory, initial $\Theta$-datum, Legendre curve,
split prime, mixed reduction, log-volume, tensor packet.

\section{Introduction}\label{s:intro}

Explicit initial $\Theta$-data are needed for any numerically effective use of the
theory; \cite{MFHMP} is devoted to obtaining such uses.  Dupuy and Hilado work out initial $\Theta$-data for the curve 11a1
over $\mathbb Q$ \cite[\S2.10]{DH25}.
Zhou describes constructions of $\mu_6$-initial $\Theta$-data from
elliptic curves over $\mathbb Q$, including Frey--Hellegouarch curves
\cite[Prop.~2.5, Lemma~2.6]{Zhou}.
These constructions concern rational moduli, whereas the mixed-place
profile considered here requires a moduli field with more than one
place above the selected rational prime.

This note contributes one,
constructed so that every condition of \cite[Def.~3.1]{IUT1} can be checked
individually, and then asks how a configuration it exhibits is handled.

The datum is stated in \S\ref{s:datum} and verified condition by condition.
\S\ref{s:dich} shows that the configuration it exhibits is not incidental, in two
senses: it is forced by the shape of \cite[Def.~3.1(e)]{IUT1}
(Proposition~\ref{p:dich}), and for this curve it cannot be removed by choosing
$\Vbad_{\mathrm{mod}}$ differently (Proposition~\ref{p:unavoid},
Remark~\ref{r:cost}).  \S\S\ref{s:sa}--\ref{s:cmp} draw a consequence, conditional
on one hypothesis about the source which we isolate and state separately;
\S\S\ref{s:fals}--\ref{s:ask} say what would falsify it and what we ask.
\S\ref{s:repro} records the computations.

Our purpose is to have the question answered.  We do not argue that the theory
fails; we were unable to locate the treatment of this configuration, and we would
regard being shown it as the useful outcome.  Two features of the situation should
be said at once.

\begin{itemize}[itemsep=2pt]
\item The configuration cannot arise when $\Fmod=\mathbb{Q}$
 (Corollary~\ref{c:dich}(i)).  This mechanism alone consequently says nothing
 about estimates for such data.  We draw no conclusion about the effective
 Diophantine results of \cite{MFHMP}.
\item That the construction below is possible at all is a property of
 \cite[Def.~3.1]{IUT1}: its conditions (a)--(f) are stated precisely enough that
 each can be verified individually against an explicit curve, several of them by
 exact machine computation.  Not every foundational definition admits this.
\end{itemize}

\section{What is and is not claimed}\label{s:scope}

We mark source statements \textup{[S]}, the present deductions and exact
calculations \textup{[P]}, and the source-adapter reading \textup{[H]}.
These are different statuses: the arithmetic does not certify \SA{}.
Statement locators for IUT I and III refer to the author's May 2020 PDFs,
IUT II to the December 2020 PDF, and IUT IV to the April 2020 PDF; the bibliography records their
2021 publication metadata separately.  For \cite{SS} we use the manuscript
dated July 16, 2018, and for \cite{Rpt} the February 2019, 45-page version
of the report on the March 2018 discussions.

\noindent\textbf{Claimed.}  (i) The datum of \S\ref{s:datum} is an initial
$\Theta$-datum with a mixed local profile above $211$.  This is independent of \SA{}; the arithmetic
in it is machine-checkable, and the remaining ingredients are standard facts and
cited constructions, identified as such in \S\ref{s:datum}.  (ii) \emph{Assuming} \SA{}, the local
comparison of \S\ref{s:cmp} fails, with margin $8/7$.

\smallskip\noindent\textbf{Not claimed.}  That the $abc$ conjecture is false;
that the lower bound of \cite[Cor.~3.12]{IUT3} is false; that the global
inequality of \cite[Thm.~1.10]{IUT4} is false; that the objections of \cite{SS}
are correct or the responses of \cite{Rpt} incorrect; that the estimate cannot be
repaired.  What is at issue is one application of one local bound.

\section{The datum}\label{s:datum}

\begin{equation}\label{e:datum}
\begin{aligned}
 &F_0=\mathbb{Q}(\sqrt5), &&\pi=16+3\sqrt5,\quad \Norm(\pi)=211,\quad a=\pi^{29},\\
 &E_0:\,y^2=x(x-1)(x-a), &&\ell=7,\\
 &F=F_0(i,E_0[15]),\quad i^2=-1, &&K=F(E_0[7])=F_0(i,E_0[105]).
\end{aligned}
\end{equation}
Let $X_F=E_F\setminus\{0\}$.  Select as $\Vbad_{\mathrm{mod}}$ all places
of $F_0$ of multiplicative reduction outside $2,7$, including the place
$v_\pi$ with $v_\pi(\pi)=1$.  The compatible section $\VV$ and cusp
$\epsilon$ are constructed below; they are not arbitrary choices.

\begin{theorem}\label{t:datum}
\textup{[P]} These data extend to an initial $\Theta$-datum in the sense of
\cite[Def.~3.1]{IUT1} and to an LGP-Gaussian log-theta-lattice.  They satisfy
the additional arithmetic hypotheses of \cite[Thm.~1.10]{IUT4}.
Above $211$, $\VV$ has exactly two elements, one selected bad and one good, and
\[
 e_{\mathrm g}=1,\qquad e_{\mathrm b}=105,\qquad
 \ord_{211}(\underline{\underline q}_{\mathrm b})=\frac{29}{7}.
\]
In particular every local index in this profile satisfies $e\le p-2=209$.
\end{theorem}

Here $\ord_p(p)=1$; an integral prime-ideal order is denoted $v_\mathfrak p$.
The following proof separates exact arithmetic from the geometric constructions.
A detailed expansion is in \cite[\S\S3--5, especially \S4.6]{Firstfull}.

\paragraph{Conditions (a) and (b).}
The field $F$ contains $i$ and full $6$-torsion, the $2$-torsion already
being rational from the Legendre equation.  Its degree over $F_0$ divides
\[
 2|\mathrm{GL}_2(\mathbb Z/15\mathbb Z)|
 =2\cdot48\cdot480=46080=2^{10}3^2 5,
\]
and it is Galois.  At every finite place, at least one of $3,5$ is prime
to the residue characteristic.  Full rational torsion of that order gives
semistable reduction by \cite[Prop.~1.8(v)]{IUT4}.  The origin supplies the
marked section; the resulting one-pointed genus-one stable model gives the
required reduction of $X_F$.

The six Legendre orbit comparisons between $a$ and $\bar a$ all fail, so
$j(E_0)\ne\overline{j(E_0)}$ and $\Fmod=\mathbb Q(j(E_0))=F_0$.
At every odd place where $a$ or $1-a$ has positive order, the integral
Legendre invariant $c_4=16(a^2-a+1)$ is a unit, proving multiplicative
reduction directly.  If $a\equiv0$, the nodal tangent equation is
$y^2=-x^2$, which splits after adjoining $i$.  If $a\equiv1$, put
$u=x-1$; the nodal tangent equation is $y^2=u^2$, already split.
Thus every selected place is split multiplicative over $F$, since $i\in F$.
At every remaining place outside $2,7$ the discriminant
$16a^2(1-a)^2$ is a unit.  Consequently the selected bad set is the full
one required for the extra good-reduction hypothesis in IV.  Also
$F_{\mathrm{tpd}}=F_0$ and $F=F_0(i,E_0[30])$ exactly, because $E_0[2]$
is already rational.  This is the exact field in IV, not an arbitrary
torsion-containing extension.

\paragraph{Condition (c), including every selected order.}
At the degree-one place $11$, $\sqrt5=4$, the reduction has trace zero;
its characteristic polynomial modulo $7$ has discriminant $5$, a
nonsquare.  Hence the representation is irreducible.  Tate inertia at
$v_\pi$, where $v_\pi(q)=58$, contains a nonidentity transvection on
$E_0[7]$.  Irreducibility provides a conjugate transvection with a different
fixed line.  Their powers generate the two opposite root groups, hence
$\mathrm{SL}_2(\mathbb F_7)$.  This group remains in the image over $F$:
the latter image is normal of index prime to $7$.  Every proper normal
subgroup of $\mathrm{SL}_2(\mathbb F_7)$ lies in its center, and the
corresponding quotient has order divisible by $7$.  Intersecting with the
latter image therefore cannot give a proper subgroup.
Thus $\mathrm{Gal}(K/F)$ has the required image.

The exact norm certificate gives
\[
 \Norm(a)=211^{29},\qquad
 \Norm(1-a)=2^2 3^2\cdot5\cdot59\cdot80621\cdot C,
\]
where
\[
 C=29622281599338016545834674091843816951690863770835007573979.
\]
The second norm is seventh-power-free, as certified by the exhaustive
sieve of \S\ref{s:repro}; primality of $C$ is not asserted.
The norm identity $v_q\Norm(z)=\sum_{\mathfrak p\mid q}
f(\mathfrak p/q)v_\mathfrak p(z)$ shows that every positive order of
$1-a$ is at most $6$.  The only positive order of $a$ is $29$ at
$v_\pi$.  Write $r$ for the relevant positive order at a selected place.
Its base Tate order is $2r$, prime to $7$.  Over a place $w$ of $F$ it
is $2r e(w/\mathfrak p)$, still prime to $7$ because the ramification
index divides $[F:F_0]$.  No selected residue characteristic is $7$.
This verifies the all-place clause, not just the order $58$ at $v_\pi$.

\paragraph{Condition (d): quotient and core.}
Fix one global line $W\subset V=E_0[7](K)$ and a nonzero
$\xi\in Q=V/W$.  Let $\pi_W:E_K\to E'=E_K/W$ and let
$\psi:E'\to E_K$ satisfy $\psi\pi_W=[7]$.  Then
$\ker\psi\simeq Q$ is constant.  Put
\[
 \underline X_K=E'\setminus\ker\psi,\qquad
 \underline C_K=[\underline X_K/\{\pm1\}],\qquad
 C_K=[X_K/\{\pm1\}].
\]
The finite etale covering $\underline C_K\to C_K$ has type
$(1,7\text{-tors})^\pm$ by \cite[Def.~2.1]{EtTh}: its geometric quotient
is onto $Q$, with zero decomposition-group action at the rational zero
cusp.  Let $\epsilon$ be the cusp corresponding to $\pm\xi$.

The four arithmetic exceptional $j$-invariants are
\[
 2^{14}31^3 5^{-3},\quad 2^2 73^3 3^{-4},\quad1728,\quad0
 \qquad\text{\cite[Prop.~2.1]{MFHMP}}.
\]
All are rational and hence exclude this $j$.  The non-arithmetic
punctured hemi-elliptic orbicurve $C_{\overline K}$ is itself a core by
\cite[Prop.~2.7]{CanLift}.  Proposition~2.3(i) there descends it to
$C_K$, the core of the finite etale cover $\underline C_K$.
For each completion, Proposition~2.3(ii) followed by (i) gives the
specified core $C_{K_v}$ as well.  Non-CM or nonintegral $j$ alone is not
being used as the criterion.  The geometric theta covers required in
(d) are supplied by \cite[Prop.~2.2, Def.~2.3]{EtTh}.

\paragraph{Conditions (e) and (f): compatible places and cusp.}
At an initially chosen lift of each selected moduli place, the split Tate
sequence $0\to\mu_7\to E[7]\to\mathbb Z/7\mathbb Z\to0$, with the
class of $q^{1/7}$ mapping to $1$, specifies a line and a nonzero quotient
class.  The group $\mathrm{SL}_2(V)$ is transitive on such marked flags:
complete a lift of the quotient class to a determinant-one basis.
Choose $\sigma\in\mathrm{Gal}(K/F)$ taking this local flag to the fixed
$(W,\pm\xi)$, and replace the chosen place by its image under $\sigma$.
The group acts on the places; the global quotient and cusp stay fixed.
Do this separately over every selected moduli place, then choose lifts
over the remaining places.  This gives a section $\VV\to\Vmod$;
equivariance of the section is not required.

Locally the quotient is $E(q)\to E(q^7)$, $z\mapsto z^7$, and its dual
is $E(q^7)\to E(q)$, $z\mapsto z$.  Thus the local cover has type
$(1,\mathbb Z/7\mathbb Z)^\pm$ and $\epsilon$ is its canonical
$\pm1$ cusp.  Full $2$-torsion gives the field condition in
\cite[Def.~3.1(e)]{IUT1}; the Weil pairings give $\mu_{12}\subset F$,
$\mu_7\subset K$.  Together with the local core, these supply the
arithmetic theta models of \cite[Thm.~1.10(iii), Def.~2.5(i)]{EtTh}.
The arrow-marked covers and their open subgroups are then those of
\cite[\S1, Def.~1.1, Rem.~1.1.2]{IUT1}, using the decomposition group of
the cusp $2\epsilon$; $2\xi\ne\pm\xi$.  The torsion assumption is on
the original $E_K$, not full $E'[7]$.  Their base changes give the covers
at good places.  This supplies the covering and cusp clauses as well as
the section.

\paragraph{Lattice and IV input.}
The native models and bridge identifications in
\cite[Prop.~4.8(ii), Prop.~6.7, Def.~6.13]{IUT1} give a full theater.
On tagged copies, the synchronized full log-links of
\cite[Prop.~1.3]{IUT3}, the nonempty enhanced Gaussian isomorphisms of
\cite[Cor.~4.10(iii)]{IUT2}, and the incoming-link Gaussian/LGP
identifications of \cite[Prop.~3.7]{IUT3} give the lattice of
\cite[Def.~3.8(ii),(iii)]{IUT3}; see the construction in
\cite[\S4.6]{Firstfull}.  This existence uses neither \SA{} nor a numerical
consequence of \cite[Cor.~3.12]{IUT3}.  The exact field and extra
good-reduction conditions of IV were checked above.

\begin{remark}[the local invariants]\label{r:loc}
At $211$, the two base places have $\sqrt5=65,146$ modulo $211$.
The corresponding residues of $a$ are $0,26$; the first has order $29$
and the second is good.  Good reduction and prime-to-$211$ torsion give
$e_{\mathrm g}=1$.  At the bad base place the nodal tangent directions
are defined by $-1$, a nonsquare in $\mathbb F_{211}$.  Over the
unramified quadratic field $k_0=\mathbb Q_{211}(i)$ the curve is split
Tate with $\ord_{211}(q)=58$.  For $n=15,105$, its exact torsion field is
$k_0(E[n])=k_0(\zeta_n,q^{1/n})$.  The denominator of $58/n$ forces its
ramification index to be a multiple of $n$; \cite[Prop.~1.8(vii)]{IUT4}
bounds that index by a divisor of $n$.  It is therefore exactly $n$.
The chosen completions of $F$ and $K$ have indices $15$ and $105$.
Finally the $2\ell$-th root satisfies
$\ord_{211}(\underline{\underline q})=58/14=29/7$.
No claim that a general root adjunction is totally ramified is needed.
\end{remark}

\section{When can this configuration occur at all?}\label{s:dich}

The following is independent of everything below; it uses only
\cite[Def.~3.1(e)]{IUT1} and \cite[Prop.~3.1]{IUT3}.

\begin{proposition}\label{p:dich}
Let $p$ be a rational prime and $\VV_p=\{v\in\VV: v\mid p\}$.  Then $|\VV_p|$
equals the number of places of $\Fmod$ above $p$.  Consequently, in the
decomposition
$\log({}^AF_{v_{\mathbb{Q}}})=\bigoplus_{\{v^\alpha\}_{\alpha\in A}}
\bigotimes_{\alpha\in A}\log({}^\alpha F_{v^\alpha})$ of \cite[Prop.~3.1]{IUT3},
the tuples $\{v^\alpha\}$ at $p$ are all constant unless $p$ has at least two
places in $\Vmod$.
\end{proposition}

\begin{proof}
\cite[Def.~3.1(e)]{IUT1} makes $\VV$ a section of
$\VV(K)\twoheadrightarrow\Vmod$, and a section restricts to a bijection above
each rational prime.  The tuple index set is $\VV_p^{\,A}$.
\end{proof}

\begin{corollary}\label{c:dich}
\textup{[P]} At a stage of length $k=j+1\ge2$, write $r=|\VV_p|$.
There are $r^k-r$ nonconstant place tuples, so they occur exactly when
$r\ge2$, independently of reduction type.  If the chosen-bad and
chosen-good subsets have sizes $b,g$, mixed-status tuples number
$r^k-b^k-g^k$ and occur exactly when $b,g>0$.  In particular:
\begin{enumerate}[label=(\roman*),itemsep=1pt]
\item if $\Fmod=\mathbb{Q}$, every tuple at every prime is constant;
\item in a quadratic moduli field, a nonsplit rational prime has a single
 place, hence only a constant tuple.  A split prime has nonconstant tuples,
 and has mixed-status tuples precisely when one place is chosen bad and
 the other chosen good.
\end{enumerate}
\end{corollary}

\begin{remark}[rational moduli data]
\label{r:Q}
By Corollary~\ref{c:dich}(i), the configuration examined in this note
\emph{cannot} arise when $\Fmod=\mathbb{Q}$.  This excludes the present
mixed-place mechanism for such data.  It does not identify every datum used
in an effective Diophantine argument with a rational-moduli datum, and we
make no claim about the validity of the effective results in \cite{MFHMP}.
\end{remark}

\begin{proposition}[the configuration cannot be removed by a choice]\label{p:unavoid}
Write $a=\pi^{29}=A+B\sqrt5$ with $A,B\in\mathbb{Z}$.  Then $\gcd(A,B)=1$ and
$\gcd(A-1,B)=3$, and consequently the only odd rational primes of multiplicative
reduction that do \emph{not} split in $F_0$ are $3$ and $5$.
\end{proposition}

\begin{proof}
Let $p$ be an odd prime inert in $F_0$, so $\sqrt5\notin\mathbb{F}_p$.  Then
$a\equiv0$ would force $A\equiv B\equiv0$, i.e.\ $p\mid\gcd(A,B)=1$, which is
impossible; and $a\equiv1$ would force $A-1\equiv B\equiv0$, i.e.\
$p\mid\gcd(A-1,B)=3$, so $p=3$ --- which is indeed inert, $5$ being a non-residue
modulo $3$.  The remaining odd non-split prime is the ramified $p=5$.
\end{proof}

\begin{remark}[what a restriction to non-split places would cost]\label{r:cost}
Write $H_{\mathrm{all}}$ and $H_{\mathrm{ns}}$ for the degree-normalized Tate
contributions outside $2\ell$ from all selected bad places and from the non-split
ones respectively.  Then
$H_{\mathrm{all}}=29\log211+\log\bigl(\Norm(1-a)/4\bigr)$,
$H_{\mathrm{ns}}=\log45$, and $H_{\mathrm{ns}}/H_{\mathrm{all}}<0.0124$, so such a
restriction would discard more than $98.76\%$ of that quantity.  (This is not a
ratio of full Weil heights.)  Independently of the cost, the additional
good-reduction hypothesis imposed in \cite[Thm.~1.10]{IUT4} does not permit an
actual bad place outside $2\ell$ to be treated as chosen-good, so the substitution
is unavailable in any case.
\end{remark}

\begin{question}\label{q:hyp}
\textup{[S]} Definition~3.1(b) of \cite{IUT1} permits a nonempty subset of
the multiplicative places of odd residue characteristic and does not impose
a splitting restriction.  Its chosen-good places may still have bad reduction.
The additional good-reduction hypothesis of \cite[Thm.~1.10]{IUT4} removes
that freedom outside $2\ell$: every actual bad place there must be selected.
Thus omitting $v_\pi$ from the present datum would prevent that IV application.

For this admitted mixed profile, how does the fixed-tuple estimate in Step~(v)
apply to the possible-image hull after the specified factor transports?
If a further restriction on the data is required, which stated hypothesis
imposes it, and which applications remain within its scope?
\end{question}

\begin{remark}[related work, added after the first deposit]\label{r:lana}
We learned of the following only after depositing the first version of this note.

\textup{[S]} \cite{LANA} isolates a compatibility problem in the passage from
\cite[Thm.~3.11]{IUT3} to \cite[Cor.~3.12]{IUT3}.  Its main goal
\cite[\S9.2]{LANA} is to show that $\eta^{q}=\eta^{\mathrm{anab}}_{S}$ for some
suitable $S$, where the two sides are isomorphisms of pointed real lines obtained
respectively from the $q$-pilot in its native structure and by anabelian and
Kummer-theoretic procedures; \cite[\S9.3]{LANA} reads this as the assertion that
the rigidified $q$-pilot is represented in the output regions.  \cite{LANA} states
that it does not have a proof of this at present, and reserves judgement.

\textup{[P]} That is not the same statement as \SA4.  \SA4 concerns which regions
the union of possible images contains, and \textup{(F2)} concerns the reading of
the resulting hull in the fixed target of Step~(v) of \cite[Thm.~1.10]{IUT4}.
Within \cite{LANA} the corresponding place is \cite[\S10.4]{LANA}, which locates
the effect of the indeterminacies in the ambiguity of the choice of $S$ and says
that this is reflected in the log-volume computation for the hull in
\cite{IUT4}.

\textup{[S]} \cite[\S10.4]{LANA} continues that the hull computation ``is made
at each place individually, but the effects of $\Ind1$, $\Ind2$ and $\Ind3$ cannot
be detected locally''.  \textup{[P]} This does not contradict
Theorem~\ref{t:cmp}, which is conditional: if that reading is correct then \SA{}
fails, and what is withdrawn is the source application, not the rational
calculation.  Nor have we identified the regions, admitted transports, hull and
normalized readout of \cite{LANA} with the fixed mixed tensor component used here;
without that identification the two statements are not yet about the same objects.
We record the cross-check as \textup{(F5)} in \S\ref{s:fals}.

\textup{[S]} One point of contact is exact.  The LGP element of
\cite[\S5.2(f)]{LANA} is
$\bigl(\bigotimes_{t\in S_j}(1)_{w\mid v_{\mathbb Q}}\bigr)\otimes
\underline{\underline q}_v^{\,j^2}$ at a bad place and $\bigl(\bigotimes_{t\in S_j}(1)_{w\mid v_{\mathbb Q}}\bigr)\otimes1$
at a good one: the factor $\underline{\underline q}_v^{\,j^2}$ occupies the distinguished slot alone and
every other slot carries $1$.  That is the same shape as the reading of Step~(v)
used in \S\ref{s:cmp}, and as the insert-$1$ representation of \S\ref{s:sa}.

\textup{[S]} \cite{LANA} gives no explicit numerical initial $\Theta$-datum
\cite[\S2]{LANA}, and computes no fibre cardinality and no separate
constant/non-constant/mixed classification of tuples.  Its containers are nevertheless indexed exactly as here:
\cite[\S5.2(a)]{LANA} sets
$K_{S_{j+1},v}=\bigl(\bigotimes_{t\in S_j}\bigoplus_{w\mid v_{\mathbb Q}}
K_{t,w}\bigr)\otimes K_{j,v}$, so the decomposition of
Proposition~\ref{p:dich} is present in that formalism and the configuration of
\S\ref{s:dich} is admitted by it.  What \S\ref{s:dich} adds is that the
configuration occurs, and \S\ref{s:datum} that it occurs in an admitted datum.

\textup{[S]} \cite{Form} decomposes the same passage as
$3.11\Rightarrow3.11.5\Rightarrow3.12$, obtained by moving the ``hull+det'' of
$3.12$ into $3.11.5$.  We use this only for the decomposition.  We do not identify
the transport of \SA3 with the algorithmic parallel transport of
\cite[Rem.~3.11.1(iv)]{IUT3}.  That remark says of the parallel transport
\emph{mechanism} that it ``does not consist of a simple instance of transport of
some set-theoretic region \dots\ via some set-theoretic function'', but is a
construction algorithm valid in the arithmetic holomorphic structures on both
sides of the $\Theta$-link.  The map $\tau_{\sigma,t}$ of \S\ref{s:sa} is of a
different type: it is a map between summands of one fixed ambient $E$ on one side.
We do not read the remark as a prohibition on such a component map, and we do not
borrow the name for it.  The label ``$3.11.5$'' is explanatory; \cite{Form} states
that it does not appear in \cite{IUT1,IUT2,IUT3,IUT4}.
\end{remark}

The selection and tuple statements preceding this question are independent of
\SA{}.  Our proposed interpretation of the comparison itself is conditional
on \SA{}, as specified next.  We do not infer an independent probability for
the splitting of bad primes from a density theorem.

\section{The hypothesis \texorpdfstring{\SA}{(SA)}}\label{s:sa}

Fix one lattice of Theorem~\ref{t:datum}, a column $(n,\circ)$,
a stage $j$ and $p=211$.  Write $E=I^{\mathbb{Q}}(S^\pm_{j+1};
{}^{n,\circ}D^\vdash_p)$ for the ambient, $U^{\mathrm{act}}_{j,t}$ for the
$t$-component of the union of possible images of a $\Theta$-pilot object in the
sense of \cite[Cor.~3.12]{IUT3}, $C_{\widetilde O_t}$ for the terminal
holomorphic hull, and $L_t=\bigotimes_i\log\mathcal{O}_{K_{t_i}}^\times$.  The source writes the tuple
decomposition explicitly: \cite[Prop.~3.1]{IUT3} gives
$\log({}^AF_{v_{\mathbb{Q}}})=\bigotimes_{\alpha\in A}\log({}^\alpha F_{v_{\mathbb{Q}}})
=\bigoplus_{\{v^\alpha\}_{\alpha\in A}}\bigotimes_{\alpha\in A}\log({}^\alpha F_{v^\alpha})$,
``the direct sum \dots\ over all collections $\{v^\alpha\}_{\alpha\in A}$'', and
\cite[Prop.~3.2]{IUT3} the mono-analytic analogue.  With $A=S^\pm_{j+1}$ the
summands are indexed by tuples $t\in\{\mathrm{g},\mathrm{b}\}^{\,j+1}$; $\Ind1$
permutes the outer label factors, while $\Ind2$ acts within the
place-indexed factors and does not freely permute the places.
Writing $\mathrm{Peel}_{\alpha,v}=\bigoplus_{t\,:\,t_\alpha=v}V_t$, a whole-factor
permutation satisfies $T_\sigma(\mathrm{Peel}_{\alpha,v})
=\mathrm{Peel}_{\sigma(\alpha),v}$: \emph{the place is unchanged and only the
position moves}.  In particular a coefficient is never re-created at the target
slot.

\begin{hypothesis}\label{h:sa}
\textup{[H]} In the fixed target arithmetic structure and normalization,
\[
\begin{cases}
 L_t\subseteq C_{\widetilde O_t}\bigl(U^{\mathrm{act}}_{j,t}\bigr),
   & \text{$t$ has at least one good factor},\\[2pt]
 p^{\lceil j^2N/7\rceil}L_t\subseteq C_{\widetilde O_t}\bigl(U^{\mathrm{act}}_{j,t}\bigr),
   & t=(\mathrm{b},\dots,\mathrm{b}),
\end{cases}
\qquad N=7\ord_{211}(\underline{\underline q}_{\mathrm{b}}).
\]
\end{hypothesis}

\SA{} is a \emph{reading} of the source, not a computation.  Its proposed
derivation has four steps.  The cited constructions are \textup{[S]};
their identification with the admitted regions and the comparison target
in these steps remains \textup{[H]}.
\begin{enumerate}[label=(SA\arabic*),leftmargin=3.1em,itemsep=1pt]
\item \emph{One output exists.}  \cite[Prop.~3.7(ii),(v)]{IUT3} makes $J_v$ ``a
 subset [indeed a submodule when $v\in\VV^{\mathrm{non}}$]'';
 \cite[Def.~3.8(i)]{IUT3} defines the $\Theta$-pilot from generators
 \emph{indexed by $v\in\Vbad$ only}; \cite[Prop.~3.10(i)]{IUT3} gives native /
 coric compatibility; \cite[Thm.~3.11(ii)(a)]{IUT3} gives isomorphisms of the
 tensor packets \emph{and} of their $\mathbb{Q}$-spans, compatible with
 log-volumes.  Our reading realizes the specified native representation as
 one region $\kappa_E[J]$ among the possible images.  It does not admit
 every isomorphic fractional ideal or replace $J$ by its full span.
\item \emph{Which automorphisms are admitted.}  In \cite[\S0, Def.~4.10]{IUT1}
 the arrows of a procession are sets of capsule-full poly-morphisms, so a factor
 transposition is realized by one automorphism of the whole procession.
\item \emph{The type of the action.}  \cite[Thm.~3.11(i)]{IUT3} treats the data
 ``regarded up to indeterminacies'', $\Ind1$ being ``the indeterminacies induced
 by the automorphisms of the procession'', and states the algorithm is
 ``functorial with respect to isomorphisms of processions''.  Labels,
 coefficients and realizations are transported together.
\item \emph{Same comparison target.}  \cite[Cor.~3.12]{IUT3} forms the
 \emph{union} of the possible images ``which we regard as subject to the
 indeterminacies $\Ind1$, $\Ind2$, $\Ind3$'', and only then the hull.
 Our reading includes the whole transported representation in this same
 union, including its global object and localizations, so that its
 $\mathcal F^\Vdash$-prime-strip restriction is preserved.  It also
 identifies the resulting fixed arithmetic target and normalized readout
 with those used in Step~(v) of \cite[Thm.~1.10]{IUT4}.
\end{enumerate}
Under these four clauses, choose $s=\sigma^{-1}t$ with a good factor in
its distinguished slot.  The native ideal generated by the insert-$1$
representation is $J_s=\widetilde O_s$ by
\cite[Prop.~3.1(ii), Prop.~3.4, Example~3.6(ii), Prop.~3.7(v), Def.~3.8(i)]{IUT3}.
Here $\widetilde O_s$ is the integral closure of the tensor order
$R_s=\bigotimes_i\mathcal O_{K_{s_i}}$, and
$L_s\subseteq R_s\subseteq J_s$ because every $e_i\le p-2$.
The ind-topological additive comparison in
\cite[Prop.~3.2, Thm.~3.11(ii)(a)]{IUT3} preserves the integral tensor
packet $I_s$ and its rational span.  On these finite-dimensional portions,
continuity gives $\mathbb Q_p$-linearity.  Since
$L_s=p^{j+1}I_s$ by \cite[Rem.~1.2.2(i)]{IUT3}, it also preserves $L$.
Transporting the whole representation, taking the $\mathbb Z_p$-span and
then the terminal target hull gives the first inclusion of \SA{}.
For the all-bad tuple, put $c=\lceil j^2N/7\rceil$ and let $a_s$ denote the
generator, at the distinguished slot of $s$, of the native local fractional ideal
$J_s$ of \cite[Prop.~3.7(v)]{IUT3} --- by \cite[Def.~3.8(i)]{IUT3} a generator up
to torsion of the monoid indexed by that place, so
$\ord(a_s)=\ord(\underline{\underline q}^{\,j^2})=j^2N/7$.  Then $p^c/a_s$ is
integral and
$p^cL_s\subseteq p^cR_s\subseteq a_sR_s\subseteq J_s$.
Integer-scalar linearity gives the second inclusion.  No ring linearity
of $\kappa_E$, or commutation of that map with the hull, is asserted.

In particular the transport has the typed component equation
\[
 \operatorname{pr}_tT_\sigma
 =\tau_{\sigma,t}\operatorname{pr}_{\sigma^{-1}t},\qquad
 \tau_{\sigma,t}:E_{\sigma^{-1}t}\longrightarrow E_t.
\]
The whole tuple is transported; no isomorphism $K_{\mathrm b}\simeq
K_{\mathrm g}$, or independent recombination of tuple components, is used.
The fuller source-to-hull argument is recorded in \cite[\S\S7--8]{Firstfull}.

\section{The comparison}\label{s:cmp}

Step (v) of \cite[Thm.~1.10]{IUT4} fixes $j$ and a collection
$\{v_i\}_{i\in S^\pm_{j+1}}$, takes ``$i^\dagger$'' to be $j\in S^\pm_{j+1}$ --- the
\emph{distinguished} slot --- and sets
\begin{quote}
``$\lambda$'' to be $0$ if $\underline v_j\in\Vgood$; ``$\lambda$'' to be
``$\ord(-)$'' of the element $\underline{\underline q}_{\,\underline v_j}^{\,j^2}$
if $\underline v_j\in\Vbad$,
\end{quote}
and then computes ``the weighted average, over possible collections
$\{v_i\}_{i\in S^\pm_{j+1}}$, of these estimates''.  So $\lambda$ is determined by
the distinguished slot alone, and the average runs over tuples.  For
$t=(\mathrm{g},\mathrm{b})$ with $j=1$ the distinguished slot is $t_1=\mathrm{b}$,
whence $\lambda=\ord_{211}(\underline{\underline q}_{\mathrm{b}}^{\,1})=29/7$; for
$j=2$ it is $4\cdot29/7=116/7$.  Note the exponent $j^2$: taking $\lambda$ to be
$\ord(\underline{\underline q})$ alone would omit it.

\begin{theorem}\label{t:cmp}
\textup{[P], conditional on [H].} Assume \SA{}.  For \eqref{e:datum} at $j=1$ and $t=(\mathrm{g},\mathrm{b})$, in
the normalization by $\log211$:
\[
 \underbrace{-\lambda+d_I+1=-\tfrac{29}{7}+\tfrac{104}{105}+1=-\tfrac{226}{105}}
   _{\text{bound from \cite[Prop.~1.2(iii)]{IUT4}}}
 \qquad\text{versus}\qquad
 \underbrace{-\tfrac{106}{105}}_{\text{hull of }L_t},
 \qquad\text{margin }\ \tfrac{8}{7}>0 .
\]
\end{theorem}

\begin{remark}[where the bound comes from, and why no error terms occur]
\label{r:tame}
\cite[Prop.~1.1]{IUT4} defines $d_i$ as the order --- for $\ord$ normalized by
$\ord(p)=1$ --- of a generator of the different of $R_i$ over $\mathbb{Z}_p$; for a
tamely ramified extension the different is $\mathfrak{p}_i^{e_i-1}$, so
$d_i=1-1/e_i$.  \cite[Prop.~1.2]{IUT4} defines $a_i$ and states: ``if $p>2$ and
$e_i\le p-2$, then $a_i=1/e_i=-b_i$''.

The bound used in Theorem~\ref{t:cmp} is \cite[Prop.~1.2(iii)]{IUT4}, stated
\emph{precisely} under the hypothesis of the present situation: ``Suppose that
$p>2$, and that $e_i\le p-2$ for all $i\in I$.  Then
$\phi(p^\lambda\cdot(R_I)^{\sim})\subseteq p^{\lambda-d_I-1}\cdot(R_I)^{\sim}$.''
Here $\lambda=29/7=435/105$ also belongs to
$e_{i^\dagger}^{-1}\mathbb Z$, as required for the chosen slot.
Every $e_i\le105\le209=p-2$ here, so it applies directly and the corresponding
normalized log-volume bound is $(-\lambda+d_I+1)\log p$, with \emph{no} error
term.  The same value follows from \cite[Prop.~1.4(iii)]{IUT4} on taking the
admissible choice $I^*=\varnothing$, when its error terms $4|I^*|/p$ and
$\sum_{i\in I^*}\{3+\log e_i\}$ both vanish; we note that the symbol $I^*$ of
\cite[Prop.~1.1]{IUT4} denotes something different.  For
$t=(\mathrm{g},\mathrm{b})$ we get $d_I=0+\tfrac{104}{105}$,
$a_I=1+\tfrac1{105}$, hence $d_I+a_I=2=|I|$, attaining with equality the
inequality $d_I+a_I\ge|I|$ of \cite[Prop.~1.4(iii)]{IUT4}.
\end{remark}

\begin{remark}[the ``nonpositive integer power'']
Enlarging the intermediate container is permitted, but the accompanying loss of
volume must then be carried into the subsequent bound.  Here
$\lfloor\lambda\rfloor-|I|=\lfloor29/7\rfloor-2=2$.
Writing $C(L_t)$ for the terminal fractional-ideal hull, hull formation
commutes with multiplication by $p^2$, so $C(p^2L_t)=p^2C(L_t)$.
This is a proper subideal of the nonzero fractional ideal $C(L_t)$.
Thus the intermediate container and its hull cannot contain the seed $L_t$;
this conclusion concerns hulls as well as the lattice inclusion.
\end{remark}

\begin{remark}[what is not refuted]
\cite[Prop.~1.4]{IUT4} itself is not in question: its estimate for the input
regions it is stated for stands.  At issue is its application to the fixed target
component of the union of possible images.
\end{remark}

\subsection*{A variant with the label $0$ fixed}\label{s:fix0}

One might ask whether fixing label $0$ in every stage eliminates the transport.
That label is distinct from the slot $j$ designated in Step~(v) when $j\ge1$.
The same datum provides a test of this proposed restriction.

\begin{theorem}\label{t:fix0}
\textup{[P], conditional on [H].} Assume \SA{} with the indicated whole
transposition still admitted.  For the same datum \eqref{e:datum}, take $j=2$, source tuple
$(\mathrm{b},\mathrm{b},\mathrm{g})$, target $t=(\mathrm{b},\mathrm{g},\mathrm{b})$
and $\sigma=(1\,2)$, which \emph{fixes the label $0$}.  Here
$\lambda=\ord(\underline{\underline q}^{\,j^2})=4\cdot\tfrac{29}{7}=\tfrac{116}{7}$,
and in the same normalization
\[
 \text{upper bound }=-\lambda+d_I+1=-\tfrac{1427}{105},
 \qquad \text{hull of }L_t=-\tfrac{107}{105},
 \qquad\text{margin }\tfrac{88}{7}>0 .
\]
\end{theorem}

Here $d_I=\tfrac{208}{105}$, $a_I=\tfrac{107}{105}$ and $d_I+a_I=3=|I|$.  In
general, writing $s=\ord_p(\underline{\underline q})$ for the depth at a bad slot,
Step (v) sets $\lambda=j^2s$ when the distinguished slot is bad.  For a tuple
with every $e_i\le p-2$ and an admitted $L_t$ seed, the margin is
\[
 \Bigl(-\sum_i e_{t_i}^{-1}\Bigr)-\bigl(-\lambda+d_I+1\bigr)=j^2s-(j+2),
\]
independently of the ramification indices, since $d_I=(j+1)-\sum_ie_{t_i}^{-1}$
in that range.  With $s=29/7$ this gives $8/7$, $88/7$ and $226/7$ for
$j=1,2,3$.  Theorems~\ref{t:cmp} and~\ref{t:fix0} are the first two cases.
Equivalently it is positive when
$\ord_p(q)>2\ell(j+2)/j^2$.  For $j=2$ this threshold is $14$.

\begin{remark}
Fixing label $0$ alone leaves the transposition $(1\,2)$ available in this
example.  This does not show that every nontrivial permutation must be
forbidden: a nontrivial subgroup fixing the actual slot $j$ can permute the
other factors.  For this seed mechanism the relevant issue is whether the
allowed orbit of $j$ reaches a good, or sufficiently shallow, source factor.
\end{remark}

\section{What would falsify this}\label{s:fals}

Theorem~\ref{t:datum} stands or falls on its arithmetic and the quoted
conditions; if a condition of \cite[Def.~3.1]{IUT1} fails, we ask that it
be identified.  The conditional numerical implication is distinct from its
proposed source application.  The latter may fail at the following points;
these are not asserted to exhaust all possible readings.
\begin{enumerate}[label=(F\arabic*),leftmargin=3.0em,itemsep=2pt]
\item The particular whole transposition reaching the good factor is not
 admitted by $\Ind1$, contrary to (SA2).  The relevant issue is the orbit
 of the actual distinguished slot, not whether some permutation is nontrivial.
\item The admitted region family or its target differs from (SA4).
 This includes failure of orbit closure for the specified native family,
 failure of its prime-strip interpretability after transport, a different
 order of span/hull and comparison, or a different identification of the
 fixed arithmetic structure and normalized readout with Step~(v).
 That step attributes the treatment of $\Ind1,\Ind2$ to the arbitrariness
 of $\phi$ in \cite[Prop.~1.2]{IUT4}.  The map there is an automorphism
 of fixed $E_t$, whereas the component transport is
 $E_{\sigma^{-1}t}\to E_t$.  A source-compatible construction connecting
 these types and explaining the applicable region and bound would answer
 the question; the type distinction alone does not rule it out.
\item The specified native representation has no admitted region with the
 claimed seed or span in the hull, contrary to (SA1).  A formal generator
 or an isomorphic ideal alone would not justify elementwise membership.
\item The comparison or subsequent transport does not carry the integral
 tensor structure and $L$ to those of the fixed target as used above.
 The log-unit equality invoked here is for $e_i\le p-2$, not for arbitrary
 tame extensions.
\item The regions, admitted transports, hull and normalized readout of
 \cite[\S9.2, \S10.4]{LANA} are identified with the fixed component used here,
 \emph{and} under that identification the effects of the indeterminacies are not
 detectable at a single place, as \cite[\S10.4]{LANA} states of the hull
 computation in \cite{IUT4}.  Both halves are needed: the identification is not
 established here, and without it the two accounts are not yet about the same
 objects.  If both are supplied then \SA{} fails and the source application is
 withdrawn.  We regard this as the cross-check most directly checkable against the
 source, since it concerns where the effect of $\Ind1$ can appear rather than our
 arithmetic.
\end{enumerate}

There is a related type distinction in
\cite[\S6.2, footnote 7]{DH}: the actual $\mathrm{ind}_1$ permits permutations between
distinct tensor factors of the adelic tensor power, which the authors'
example does not exhibit.  They then derive an upper bound of
their own in \cite[\S6.3, Lemma~6.3.1]{DH}.  We cite this observation for that distinction only, not as
verification of \SA{} or of the present conclusion.

\textbf{(F2)} specifies which source quantity is compared.  We do not
identify it with the separate objection answered in \cite[\S12]{Rpt}, or
with the argument in \cite{SS}; those comparisons require their own
source and normalization analysis.  A demonstration that our source
application fails at any of these points would settle that application.
\emph{We would withdraw the source application of
Theorems~\ref{t:cmp} and~\ref{t:fix0}, and would regard the demonstration
as the more valuable outcome}.  The conditional rational calculation
would remain, with its unsuccessful source-adapter premise identified.

\section{What we ask}\label{s:ask}

First, we ask that the admitted initial datum in \S\ref{s:datum} and its
mixed profile be checked against the source conditions.  Its construction
does not depend on \SA{}, and dropping its split bad place outside $2\ell$
does not preserve the extra hypothesis of IV.  The nonconstant and
mixed-status tuple statements in \S\ref{s:dich} are distinct elementary
consequences of the section defining $\VV$.

Second, we ask whether the full reading (SA1)--(SA4), including the
fixed target and its normalized readout, is correct for this datum.
Theorems~\ref{t:cmp} and~\ref{t:fix0} are offered as conditional local
comparisons and a request to decide that reading, not as a settled
adjudication of the theory.  If the application is wrong, identifying the
failed clause and the source-compatible comparison would answer the question.

\section{Reproduction}\label{s:repro}

The companion package \cite{Firstfull} separates the checks from the source
arguments.  The standard-library program \texttt{scripts/checks.py}
reproduces the quadratic-integer arithmetic, Frobenius traces, rational
local constants, and corrected stage margins.  Its profile identities
do not prove the existence of arbitrary initial data.  The additional
programs \path{scripts/check_round3.py} and
\path{scripts/check_normalization.py} record the later exact checks
and the separate normalization identities.

The all-selected-order certificate requires more than those short checks:
\path{scripts/check_norm_seventh_power_free.cpp} and its log test all
primes up to $\lfloor C^{1/7}\rfloor=225469788$.  The norm identity then
translates that integer certificate to the selected prime-ideal orders.
This exhaustive sieve does not certify primality of $C$.
The geometric constructions, cited theorems, and \SA{} are not verified
by these programs.  In particular the existence of the lattice and the
compatible marked covers is the mathematical construction in
\S\ref{s:datum}, expanded in \cite[\S4]{Firstfull}.

The distinction between invariance of individual log-volumes and stability
of a family of possible images is also illustrated by a finite model
checked by the Lean~4 kernel in the companion package.  Even both
properties do not bound a union hull by a member hull.  This model is
not a formalization of the source-adapter hypothesis or an IUT counterexample.

\section*{Disclosure and corrections}

Large language models were used substantially in locating passages,
cross-checking computations, drafting, and adversarial review.
The author is responsible for the claims.  The exact computations and the
source arguments are separately inspectable; \SA{} remains a proposed
reading of the source, with its failure conditions stated above.

The correction history is material.  An earlier version used $29/7$
instead of $j^2(29/7)$ at $j=2$, giving the erroneous margin $1/7$.
The correct upper bound is $-1427/105$ and the margin is $88/7$.
The old value is not a second instance of the same datum.
Earlier explanations also wrongly used full $2$-torsion to exclude a
quadratic twist, and used only the singleton bad set $\{v_\pi\}$ despite
the extra IV condition.  Both have been corrected.
The power of $2$ in the degree bound is $2^{10}$, not $2^{11}$.

A conjecture that all four arithmetic exceptional $j$-invariants were
algebraic integers was false for two of them.  The proof uses their
rationality and the nonrationality of the present $j$, with the
factorization $488095744=2^{14}31^3$ checked against
\cite[Prop.~2.1]{MFHMP}.  We have also corrected the conflation of
nonconstant tuples with mixed reduction, the generalization of the
log-unit equality beyond $e\le p-2$, and the claim that fixing label $0$
would require banning every nontrivial permutation.  We do not claim
that no further errors remain, and retain the withdrawal undertaking of
\S\ref{s:fals}.

\begin{thebibliography}{MFHMP}
\bibitem[First-full]{Firstfull} T.~Urahigashi, \emph{An explicit initial $\Theta$-datum
 with mixed reduction over a split prime, and the local estimates of
 inter-universal Teichm\"uller theory IV}, companion draft, version
 0.1.3-draft (2026).  Included as \href{paper.pdf}{\texttt{paper.pdf}},
 with the source and reproducibility scripts in the same package.
\bibitem[CanLift]{CanLift} S.~Mochizuki, \emph{The absolute anabelian geometry of
 canonical curves}, Doc. Math., Kazuya Kato's fiftieth birthday, Extra Vol.\
 (2003), 609--640.
\bibitem[DH]{DH} T.~Dupuy, A.~Hilado, \emph{Probabilistic Szpiro, Baby Szpiro,
 and Explicit Szpiro from Mochizuki's Corollary 3.12}, arXiv:2004.13108v2
 [math.NT] (2020).
\bibitem[EtTh]{EtTh} S.~Mochizuki, \emph{The \'etale theta function and its
 Frobenioid-theoretic manifestations}, Publ. Res. Inst. Math. Sci.
 \textbf{45} (2009), no.~1, 227--349.
\bibitem[IUT1]{IUT1} S.~Mochizuki, \emph{Inter-universal Teichm\"uller theory I},
 Publ. Res. Inst. Math. Sci. \textbf{57} (2021), 3--207.
 Consulted author version: May 2020.
\bibitem[IUT2]{IUT2} S.~Mochizuki, \emph{Inter-universal Teichm\"uller theory II},
 Publ. Res. Inst. Math. Sci. \textbf{57} (2021), 209--401.
 Consulted author version: December 2020.
\bibitem[IUT3]{IUT3} S.~Mochizuki, \emph{Inter-universal Teichm\"uller theory III},
 Publ. Res. Inst. Math. Sci. \textbf{57} (2021), 403--626.
 Consulted author version: May 2020.
\bibitem[IUT4]{IUT4} S.~Mochizuki, \emph{Inter-universal Teichm\"uller theory IV},
 Publ. Res. Inst. Math. Sci. \textbf{57} (2021), 627--723.
 Consulted author version: April 2020.
\bibitem[Form]{Form} S.~Mochizuki, \emph{On the formalization of IUT: a preliminary
 progress report} (joint work in progress with Y.~Hoshi, G.~Yamashita, Y.~Yang),
 slides, Workshop on AI and Theorem Provers in Mathematics, April 2026, 14 pp.
\bibitem[LANA]{LANA} Project LANA, \emph{Project LANA interim report on IUT theory},
 July 2026, 50 pp.  Lean for ANAbelian geometry; core members J.~Commelin,
 K.~Kedlaya, Y.~Hoshi, A.~Topaz, F.~Kato (project leader).
\bibitem[MFHMP]{MFHMP} S.~Mochizuki, I.~Fesenko, Y.~Hoshi, A.~Minamide,
 W.~Porowski, \emph{Explicit estimates in inter-universal Teichm\"uller theory},
 Kodai Math. J. \textbf{45} (2022), no.~2, 175--236.
 Consulted author versions: December 2021 and May 2022.
\bibitem[Rpt]{Rpt} S.~Mochizuki, \emph{Report on discussions, held during the
 period March 15--20, 2018, concerning inter-universal Teichm\"uller theory},
 RIMS, Kyoto University; consulted February 2019 version, 45 pp.
\bibitem[SS]{SS} P.~Scholze, J.~Stix, \emph{Why $abc$ is still a conjecture},
 manuscript, July 16, 2018.
\bibitem[DH25]{DH25} T.~Dupuy, A.~Hilado,
\emph{The statement of Mochizuki's Corollary 3.12: initial theta data},
arXiv:2004.13228v2 [math.NT], 19 June 2025, 20~pp.
\bibitem[Zhou]{Zhou} Z.-P.~Zhou,
\emph{The inter-universal Teichm\"uller theory and new Diophantine
results over the rational numbers.~I},
arXiv:2503.14510v1 [math.NT], submitted 8 March 2025, 41~pp.
\end{thebibliography}

\end{document}
